\chapter{Conclusiones}
\label{ch:conclusiones}

Tras el análisis e implementación del Ejercicio Práctico N° 1, se concluye que se han alcanzado satisfactoriamente
los objetivos propuestos. La utilización de una interrupción externa gestionada a través del NVIC y el módulo EXTI
ha permitido resolver de manera eficiente la detección y el procesamiento de eventos asíncronos (la pulsación del
botón en PC15), rompiendo el flujo secuencial del programa principal sin requerir un muestreo periódico (polling).

El desarrollo en un entorno Bare Metal puro (sin sistemas operativos ni capas de abstracción complejas) ha resultado
fundamental para comprender de manera profunda el comportamiento a bajo nivel de la arquitectura ARM Cortex-M3. La
configuración manual de registros específicos como los de habilitación de reloj en el RCC, control de puertos GPIO,
asociación de líneas en el AFIO, configuración de flancos en el EXTI, y activación en el NVIC proporciona un control
absoluto sobre el hardware y optimiza el uso de recursos de memoria y CPU.

Asimismo, la implementación de la tabla de vectores en el archivo de inicio en ensamblador y la definición exacta de la
distribución física en el script de enlazado ponen de manifiesto la estrecha relación que existe entre el software
desarrollado y la estructura del hardware subyacente. Se constata la vital importancia de limpiar el flag pendiente
de interrupción en el registro correspondiente antes de finalizar la rutina de servicio, ya que la omisión de este
paso invalida el comportamiento esperado del microcontrolador al generar bucles infinitos de atención.

